% English translation of checked_pages/esq_119_017.tex; source: cote 119, site 17, PDF page 18, typescript page 15.
% AI-translated by GPT-6 Astra (OpenAI, ChatGPT Pro mode).
\scanpage{119/17}
\mspage{15}

embarrassed to submit it to the experts in the subject. Deligne, consulted,
did indeed find the supposition crazy\frapparatus{V-FR-017-001}, but had no counterexample
up his sleeve. Less than a year later, at the International
Congress in Helsinki, the Soviet mathematician Bielyi\tnote{TN-017-01}{The source's spelling ``Bielyi'' is retained.}
announced precisely this result, with a proof of disconcerting simplicity
occupying two short pages of a letter from Deligne - probably never
has a deep and bewildering result been proved
in so few lines !

In the form in which Bielyi states it, his result says essentially
that \uline{every algebraic curve defined over a number}
\uline{field can be obtained as a covering of the projective line ramified}
\uline{only at the points} 0, 1, $\infty$. This result seems to have passed
more or less unnoticed. Yet it appears to me to have considerable
scope. For me, its essential message was that \uline{there is a deep}
\uline{identity between the combinatorics of finite maps on the one hand}, \uline{and the}
\uline{geometry of algebraic curves defined over number fields}
\uline{on the other}. This deep result, together with the algebraic-geometric
interpretation of finite maps, opens the door onto a new,
unexplored world - and within reach of everyone, who passes by without seeing it.

It was only almost three years later, seeing that the vast
horizons opening up here decidedly stirred
nothing in any of my students, or even in any of the three or
four colleagues of the highest caliber to whom I had the opportunity
to speak about it in detail, that I made a first exploratory
journey into this ``new world'', from January to June 1981. This first venture
took material form in a packet of some 1 300 manuscript pages,
christened ``The Long March through Galois Theory''.\tnote{TN-017-02}{The French title is ``La Longue Marche à travers la théorie de Galois''.}
It is above all an effort to understand the relations between
``arithmetic'' Galois groups and ``geometric'' profinite fundamental
groups. Fairly quickly it turned toward formulating
the action of $\mathrm{Gal}(\bar{\mathbb{Q}}/\mathbb{Q})$ on $\hat{\pi}_{0,3}$ in calculational terms, and at a
later stage, on the slightly larger group $S\ell(2,\mathbb{Z})^{\wedge}$, which
gives rise to a more elegant and more effective formalism. It was also
in the course of this work (but developed in separate notes)

\margin{17}

